\section{Szeroka całka Denjoya (Denjoya--Chinczyna)}
Zmieniamy dwa elementy definicji: oscylacje zastępujemy przyrostami między końcami przedziałów, a zwykłą pochodną --- pochodną aproksymatywną. Zachowujemy ciągłość funkcji pierwotnej.
\begin{definicja}[Pochodna aproksymatywna]
Mierzalny zbiór $E\subset\R$ ma gęstość jeden w punkcie $x$, jeśli $\lambda_1(E\cap(x-r,x+r))/(2r)\to1$ przy $r\downarrow0$. Gęstość zero definiujemy analogicznie, zastępując granicę przez zero. Pochodną aproksymatywną rozpatrujemy w punktach wnętrza przedziału; wartości na końcach nie wpływają na całkę. Mówimy, że $F$ ma pochodną aproksymatywną $F'_{\mathrm{ap}}(x)=L$, jeżeli dla każdego $\varepsilon>0$ zbiór punktów $y\ne x$ spełniających
$\left|\frac{F(y)-F(x)}{y-x}-L\right|\ge\varepsilon$ ma gęstość zero w $x$. Zwykła pochodna, o ile istnieje, jest też pochodną aproksymatywną o tej samej wartości.
\end{definicja}
\begin{lemat}[Jednoznaczność i liniowość pochodnej aproksymatywnej]\label{it:lem:approx-derivative}
Skończona pochodna aproksymatywna w punkcie jest wyznaczona jednoznacznie. Jeśli dwie funkcje mają w tym punkcie skończone pochodne aproksymatywne, to ich kombinacja liniowa ma pochodną aproksymatywną równą odpowiedniej kombinacji.
\end{lemat}
\begin{proof}
Gdyby dwie różne liczby $L_1,L_2$ były pochodnymi aproksymatywnymi tej samej funkcji, bierzemy $\varepsilon=|L_1-L_2|/3$. Zbiory punktów, w których iloraz różnicowy różni się od $L_i$ o mniej niż $\varepsilon$, mają gęstość jeden. Ich przecięcie również ma gęstość jeden, ponieważ miara dopełnienia przecięcia jest ograniczona przez sumę miar dopełnień. Przecięcie nie jest więc puste w dostatecznie małych otoczeniach punktu, ale nierówność trójkąta dawałaby $|L_1-L_2|<2\varepsilon$, sprzeczność.
Dla kombinacji $\alpha F+\beta G$ wybieramy dla obu funkcji zbiory dobrych ilorazów z dokładnością $\varepsilon/(1+|\alpha|+|\beta|)$. Na ich przecięciu błąd ilorazu kombinacji jest mniejszy od $\varepsilon$, a przecięcie ma gęstość jeden. To daje tezę.
\end{proof}

\begin{definicja}[Klasa $\AC G$ i szeroka całka Denjoya]
Dla ciągłej $F:[a,b]\to\R$ i $E\subset[a,b]$ piszemy $F\in\AC(E)$, jeśli dla każdego $\varepsilon>0$ istnieje $\eta>0$ takie, że dla każdej skończonej rodziny przedziałów $[u_j,v_j]$ o rozłącznych wnętrzach i końcach należących do $E$ warunek $\sum_j(v_j-u_j)<\eta$ pociąga $\sum_j|F(v_j)-F(u_j)|<\varepsilon$. Piszemy $F\in\AC G[a,b]$, gdy $[a,b]=\bigcup_{n\ge1}E_n$ dla zbiorów domkniętych $E_n$ z $F\in\AC(E_n)$. Funkcja $f$ jest całkowalna w szerokim sensie Denjoya, jeżeli istnieje ciągła $F\in\AC G[a,b]$ z $F'_{\mathrm{ap}}=f$ prawie wszędzie. Definiujemy $ (\mathrm D_{\mathrm{sz}})\int_a^b f=F(b)-F(a)$.
\end{definicja}
\begin{lemat}[Przedłużenie ze zbioru domkniętego]\label{it:lem:ac-affine-extension}
Jeśli $E\subset[a,b]$ jest domknięty, $F:[a,b]\to\R$ jest ciągła i $F\in\AC(E)$, to $F|_E$ ma absolutnie ciągłe przedłużenie $u:[a,b]\to\R$. Ponadto dla każdego zbioru borelowskiego $B\subset E$ zachodzi
\[
 \lambda_1^*(F(B))\le\int_B|u'(x)|\,\mathrm d\lambda_1(x).
\]
W szczególności dla $B=E$ obraz jest zwarty, więc można użyć miary zamiast miary zewnętrznej.
\end{lemat}
\begin{proof}
Dla pustego lub jednopunktowego $E$ bierzemy przedłużenie stałe. W pozostałych przypadkach oznaczamy luki wewnętrzne przez $(l_k,r_k)$ i przedłużamy $F$ afinicznie na ich domknięciach, a poza $[\min E,\max E]$ stale. Ciągłość tego przedłużenia wynika z argumentu o krótkich i długich lukach z lematu \ref{it:lem:acstar-affine-extension}; ten argument nie wymaga kontroli oscylacji.

Najpierw wykażemy, że $\sum_k|F(r_k)-F(l_k)|<\infty$. Wybieramy $\eta_0$ z warunku $\AC(E)$ dla błędu $1$. Ponieważ suma długości luk nie przekracza $b-a$, po odrzuceniu skończenie wielu luk suma pozostałych długości jest mniejsza od $\eta_0$. Dla każdej skończonej podrodziny ogona warunek $\AC(E)$ ogranicza sumę modułów przyrostów przez $1$. Supremum po podrodzinach daje zbieżność szeregu.
Stąd funkcja
\[
 w(x)=\sum_k\frac{|F(r_k)-F(l_k)|}{r_k-l_k}\chi_{(l_k,r_k)}(x)
\]
należy do $\mathcal L^1$: jest nieujemna, mierzalna, a jej całka na mocy zbieżności monotonicznej jest powyższym szeregiem.

Dla przedziału $[s,t]$ przecinającego $E$ połóżmy
$p=\min(E\cap[s,t])$ i $q=\max(E\cap[s,t])$. Rozpisując przyrost przez $p,q$, otrzymujemy
\[
 |u(t)-u(s)|\le |F(q)-F(p)|+
 \int_{[s,p]\cup[q,t]}w\,\mathrm d\lambda_1.
\]
Gdy przedział nie przecina $E$, moduł przyrostu jest równy całce $w$ po nim (również w luce brzegowej, gdzie oba są zerowe).
Dla skończonej rodziny przedziałów o rozłącznych wnętrzach środkowe odcinki $[p,q]$ mają końce w $E$, rozłączne wnętrza i sumę długości nie większą od długości całej rodziny. Przy dostatecznie małej tej sumie ich przyrosty są ograniczone przez $\varepsilon/2$ z warunku $\AC(E)$. Wszystkie pozostałe odcinki mają rozłączne wnętrza i łączną długość nie większą od długości rodziny. Absolutna ciągłość całki $w$ daje dla dostatecznie małej długości ich łączny wkład mniejszy od $\varepsilon/2$. Dowodzi to $u\in\AC[a,b]$.

Z twierdzenia \ref{it:thm:ac-lebesgue-primitive} wynika $\operatorname{diam}u(J)\le\int_J|u'|\,\mathrm d\lambda_1$ dla każdego przedziału $J\subset[a,b]$ (także względnie otwartego, przez supremum po jego punktach).
Dla $\varepsilon>0$ z regularności miary i absolutnej ciągłości całki wybieramy względnie otwarty $O\supset B$ z $\int_{O\setminus B}|u'|\,\mathrm d\lambda_1<\varepsilon$. Rozkładamy $O$ na co najwyżej przeliczalnie wiele rozłącznych przedziałów $J_k$. Obrazy $u(J_k)$ są przedziałami z ciągłości $u$, a ich długości nie przekraczają całek $|u'|$ po $J_k$. Ponieważ $F(B)=u(B)\subset\bigcup_k u(J_k)$, podaddytywność miary zewnętrznej daje
\[
 \lambda_1^*(F(B))\le\sum_k\int_{J_k}|u'|\,\mathrm d\lambda_1
 =\int_O|u'|\,\mathrm d\lambda_1
 <\int_B|u'|\,\mathrm d\lambda_1+\varepsilon.
\]
Dowolność $\varepsilon$ kończy dowód.
\end{proof}
\begin{twierdzenie}[Pochodna aproksymatywna i jednoznaczność]\label{it:thm:acg-uniqueness}
Każda ciągła funkcja $F\in\AC G[a,b]$ ma skończoną pochodną aproksymatywną prawie wszędzie. Jeśli $F'_{\mathrm{ap}}=0$ prawie wszędzie, to $F$ jest stała. W konsekwencji szeroka całka Denjoya jest dobrze określona i liniowa.
\end{twierdzenie}
\begin{proof}
Niech $[a,b]=\bigcup_n E_n$ będzie domkniętym pokryciem z definicji, a $u_n$ --- absolutnie ciągłym przedłużeniem $F|_{E_n}$ z lematu \ref{it:lem:ac-affine-extension}. Prawie każdy punkt $x\in E_n$ jest punktem gęstości jeden zbioru $E_n$ i punktem różniczkowalności $u_n$. W takim punkcie, dla $y\in E_n$, ilorazy różnicowe $F$ są równe ilorazom dla $u_n$, więc $F'_{\mathrm{ap}}(x)=u_n'(x)$. Wyjątki ze wszystkich $n$ tworzą zbiór miary zero. Pochodna aproksymatywna ma mierzalnego reprezentanta: na zbiorze $E_n\setminus\bigcup_{j<n}E_j$ bierzemy mierzalną wersję $u_n'$ z zerem na zbiorze nieróżniczkowalności. Te zbiory są borelowskie i tworzą rozłączne pokrycie. Otrzymana funkcja zgadza się z pochodną aproksymatywną prawie wszędzie, zatem każda jej wersja o skończonych wartościach na zbiorze wyjątkowym jest mierzalna w sensie Lebesgue'a.

Jeżeli $F'_{\mathrm{ap}}=0$ prawie wszędzie, to $u_n'=0$ prawie wszędzie na $E_n$. Z lematu wynika $\lambda(F(E_n))\le\int_{E_n}|u_n'|\,\mathrm d\lambda=0$, a więc $\lambda(F([a,b]))=0$. Ponieważ ciągły obraz przedziału jest przedziałem, $F([a,b])$ musi być jednym punktem i $F$ jest stała.

Klasa $\AC G$ jest stabilna na kombinacje liniowe: wystarczy przejść do przeliczalnego pokrycia przecięciami zbiorów domkniętych. Liniowość pochodnych aproksymatywnych wynika z lematu \ref{it:lem:approx-derivative}. Dwie pierwotne tej samej funkcji różnią się zatem stałą, a całka jest jednoznaczna i liniowa.
\end{proof}
\begin{wniosek}[Właściwe rozszerzenie klasy HK]
Każda funkcja całkowalna w sensie HK, Perrona lub wąskim sensie Denjoya jest całkowalna w szerokim sensie Denjoya z tą samą wartością całki. Inkluzja jest właściwa.
\end{wniosek}
\begin{proof}
Dla dowolnego $E$ nierówność $|F(v)-F(u)|\le\osc_{[u,v]}F$ daje $\AC G_*\subset\AC G$, a zwykła pochodna jest pochodną aproksymatywną. Właściwość wykaże następny przykład.
\end{proof}
\begin{przyklad}[Oscylacje w lukach zbioru Cantora]
Niech $K\subset[0,1]$ będzie klasycznym zbiorem Cantora. Na etapie $n\ge1$ usuwanych jest $2^{n-1}$ luk $I_{n,k}=(a_{n,k},b_{n,k})$ długości $3^{-n}$. Definiujemy
\[
F(x)=\begin{cases}
2^{-n}\sin^2\!\bigl(\pi3^n(x-a_{n,k})\bigr),&x\in I_{n,k},\\
0,&x\in K.
\end{cases}
\]
Funkcja jest ciągła: na końcach luk wynosi zero, a amplitudy na lukach wysokich poziomów dążą do zera. Jest klasy $\AC$ na domknięciu każdej luki i stała na $K$, zatem należy do $\AC G[0,1]$.

Pokażemy, że $F\notin\AC G_*[0,1]$. Gdyby $K$ było pokryte przeliczalnie wieloma domkniętymi zbiorami $E_m\cap K$ z $F\in\AC_*(E_m)$, twierdzenie Baire'a zastosowane w $K$ dałoby podstawowy przedział Cantora $J$ pewnego poziomu $r$, dla którego $J\cap K\subset E_m$ dla jednego $m$. Wszystkie luki poziomu $n>r$ zawarte w $J$ mają końce w $E_m$. Ich liczba wynosi $2^{n-r-1}$, suma długości wynosi $2^{n-r-1}3^{-n}\to0$, a suma oscylacji jest równa $2^{n-r-1}2^{-n}=2^{-r-1}$, co przeczy definicji $\AC_*(E_m)$.

Niech $f=F'$ na dopełnieniu $K$, a na $K$ określmy $f=0$. W lukach mamy
$f(x)=\pi(3/2)^n\sin(2\pi3^n(x-a_{n,k}))$.
Ponieważ $\lambda(K)=0$, $F'_{\mathrm{ap}}=f$ prawie wszędzie. Stąd $(\mathrm D_{\mathrm{sz}})\int_0^1 f=F(1)-F(0)=0$. Gdyby $f$ była całkowalna HK, jej pierwotna $G$ należałaby do $\AC G_*$. Różnica $G-F$ miałaby zerową pochodną aproksymatywną prawie wszędzie, więc na mocy poprzedniego twierdzenia byłaby stała, wbrew $F\notin\AC G_*$. Ponadto na każdej luce poziomu $n$ całka modułu pochodnej wynosi $2\cdot2^{-n}$: funkcja rośnie od zera do $2^{-n}$ i maleje do zera. Na całym poziomie suma wynosi $2^{n-1}\cdot2^{1-n}=1$. Z addytywności całki nieujemnej po rozłącznych lukach otrzymujemy $\int_0^1|f|\,\mathrm d\lambda_1=\sum_{n\ge1}1=\infty$, więc $f\notin\mathcal L^1$.
\end{przyklad}
\begin{uwaga}[Granica teorii]
Szeroka całka Denjoya odtwarza funkcje pierwotne klasy $\AC G$ z ich pochodnych aproksymatywnych, lecz nie dowolne funkcje ciągłe z pochodnych prawie wszędzie. Funkcja Cantora ma pochodną aproksymatywną zero prawie wszędzie, ale nie jest klasy $\AC G$, gdyż jest niestała.
\end{uwaga}

% ============================================================
% UZUPELNIENIE MATEMATYCZNE: ROZDZIAL 9
% ============================================================
\subsection*{Różnica między $\AC(E)$ i $\AC_*(E)$}
\begin{przyklad}[Brak kontroli oscylacji na lukach]\label{prz:ac-kontra-acstar}
Niech $K$ będzie klasycznym zbiorem Cantora, a $F$ funkcją z przykładu o oscylacjach w lukach $K$ z tego rozdziału. Na $K$ zachodzi $F=0$, więc $F\in\AC(K)$: przyrosty pomiędzy dowolnymi punktami $K$ są równe zeru. Nie zachodzi natomiast $F\in\AC_*(K)$. Istotnie, dla luk $I_{n,k}$ usuwanych na poziomie $n$ mamy końce w $K$ oraz
\[
\sum_{k=1}^{2^{n-1}}|I_{n,k}|=2^{n-1}3^{-n}\longrightarrow0,
\qquad
\sum_{k=1}^{2^{n-1}}\osc_{\overline I_{n,k}}F
=2^{n-1}2^{-n}=\frac12.
\]
Dla każdego $\eta>0$ znajdziemy zatem rozłączne przedziały o końcach w $K$ i łącznej długości mniejszej niż $\eta$, na których suma oscylacji pozostaje równa $1/2$. To pokazuje, dlaczego w szerokiej teorii kontrolujemy jedynie przyrosty, a w wąskiej całe oscylacje.
\end{przyklad}
\begin{twierdzenie}[Wzór Newtona--Leibniza dla szerokiej całki Denjoya]\label{tw:denjoy-szeroka-ftc}
Niech $F\in C[a,b]\cap\AC G[a,b]$. Wówczas $F'_{\mathrm{ap}}$ istnieje i jest skończona prawie wszędzie. Ustalmy dowolną funkcję $f:[a,b]\to\R$ o skończonych wartościach, równą $F'_{\mathrm{ap}}$ wszędzie tam, gdzie ta pochodna istnieje poza pewnym zbiorem miary zero. Wówczas dla dowolnych $a\le u\le v\le b$ funkcja $f$ jest całkowalna w szerokim sensie Denjoya na $[u,v]$ i
\[
 (\mathrm D_{\mathrm{sz}})\int_u^v f(x)\,\mathrm dx=F(v)-F(u).
\]
\end{twierdzenie}
\begin{proof}
Istnienie pochodnej aproksymatywnej prawie wszędzie ustalono we wcześniejszym twierdzeniu tego rozdziału. Jeśli $[a,b]=\bigcup_n E_n$ jest domkniętym pokryciem z $F\in\AC(E_n)$, to zbiory $E_n\cap[u,v]$ są domknięte w $[u,v]$ i pokrywają ten przedział; dlatego $F|_{[u,v]}\in\AC G[u,v]$. W punktach wewnętrznych $[u,v]$ pochodna aproksymatywna ograniczenia zgadza się z pochodną aproksymatywną $F$, o ile istnieje. Definicja szerokiej całki Denjoya daje tożsamość. Jej jednoznaczność wynika z wcześniejszego twierdzenia o zerowej pochodnej aproksymatywnej.
\end{proof}
\begin{stwierdzenie}[Lokalność i liniowość szerokiej całki]\label{st:denjoy-szeroka-lokalnosc}
Jeśli $f,g$ są całkowalne w szerokim sensie Denjoya na $[a,b]$ i $\alpha,\beta\in\R$, to $\alpha f+\beta g$ jest całkowalna i jej całka jest odpowiednią kombinacją liniową. Jeśli $a<c<b$, to $f$ jest całkowalna na $[a,b]$ wtedy i tylko wtedy, gdy jest całkowalna na $[a,c]$ i $[c,b]$; wówczas całki są addytywne.
\end{stwierdzenie}
\begin{proof}
Dla pierwotnych $F,G$ bierzemy $H=\alpha F+\beta G$. Przez przecięcia domkniętych zbiorów pokrycia uzyskujemy $H\in\AC G$, a pochodna aproksymatywna jest liniowa wszędzie, gdzie obie pochodne istnieją. Dla lokalności ograniczamy pierwotną do podprzedziałów. Odwrotnie, dwie pierwotne na podprzedziałach można skleić, dodając jednej z nich stałą, aby wartości w $c$ się zgadzały. Otrzymujemy ciągłą funkcję, a połączone domknięte pokrycia obu przedziałów świadczą, że jest klasy $\AC G$. Przyrost tej funkcji jest sumą przyrostów po obu częściach.
\end{proof}
\begin{wniosek}[Własność Luzina i ograniczone wahanie]\label{it:cor:acg-bv}
Każda funkcja $F\in\AC G[a,b]$ ma własność Luzina (N): obraz zbioru miary Lebesgue'a zero ma miarę zero. Ponadto
\[
 \AC G[a,b]\cap\BV[a,b]=\AC[a,b].
\]
Jeśli $f$ jest całkowalna w szerokim sensie Denjoya, a $F(x)=(\mathrm D_{\mathrm{sz}})\int_a^x f$, to $F\in\BV[a,b]$ wtedy i tylko wtedy, gdy $f\in\mathcal L^1[a,b]$. Wówczas $V_a^b(F)=\int_a^b|f|\,\mathrm d\lambda_1$.
\end{wniosek}
\begin{proof}
Niech $E_n$ będzie domkniętym pokryciem z definicji $\AC G$, a $u_n$ absolutnie ciągłymi przedłużeniami z lematu \ref{it:lem:ac-affine-extension}. Dla zbioru zerowego $N$ wybieramy borelowski zbiór zerowy $B\supset N$. Każdy $B\cap E_n$ jest borelowski, zatem
$\lambda_1^*(F(B\cap E_n))\le\int_{B\cap E_n}|u_n'|\,\mathrm d\lambda_1=0$.
Z przeliczalnej podaddytywności miary zewnętrznej wynika
$\lambda_1^*(F(N))=0$.

Jeśli $F\in\AC G\cap\BV$, z rozkładu Jordana i twierdzenia o różniczkowaniu funkcji niemalejących (argument z lematu \ref{it:lem:bvstar-null-to-acstar}) wynika, że zwykła pochodna $F'$ istnieje prawie wszędzie i jest całkowalna w sensie Lebesgue'a po uzupełnieniu na zbiorze wyjątkowym. Zwykła pochodna jest również aproksymatywna. Funkcja
\[
 G(x)=F(a)+\int_a^x F'(t)\,\mathrm d\lambda_1(t)
\]
jest absolutnie ciągła, a więc należy do $\AC G$, i ma zwykłą pochodną $F'$ prawie wszędzie (twierdzenie \ref{it:thm:ac-lebesgue-primitive}). Stąd $F-G$ należy do $\AC G$ i ma zerową pochodną aproksymatywną prawie wszędzie. Twierdzenie \ref{it:thm:acg-uniqueness} daje $F-G$ stałą; wartość w $a$ jest zerowa, więc $F=G\in\AC$. Inkluzja odwrotna wynika z definicji oraz ograniczonego wahania funkcji absolutnie ciągłej.

Dla pierwotnej szerokiej całki Denjoya mamy $F'_{\mathrm{ap}}=f$ prawie wszędzie. Jeśli $F\in\BV$, pierwsza część daje $F\in\AC$, a zwykła pochodna $F'$ jest równa $f$ prawie wszędzie przez jednoznaczność pochodnej aproksymatywnej (lemat \ref{it:lem:approx-derivative}). Zatem $f\in\mathcal L^1$. Odwrotnie, jeśli $f\in\mathcal L^1$, jego pierwotna Lebesgue'a jest absolutnie ciągła i ma pochodną aproksymatywną równą $f$ prawie wszędzie. Jednoznaczność pierwotnych szerokiej całki Denjoya daje jej równość z $F$. Wzór na wahanie wynika teraz z twierdzenia \ref{it:thm:denjoy-bv-ac}, ponieważ $F$ jest absolutnie ciągła, a więc jest także pierwotną wąskiej całki Denjoya.
\end{proof}

\begin{uwaga}[Nazewnictwo i oddzielne poziomy teorii]
W terminologii historycznej szeroka całka Denjoya jest także nazywana całką Denjoya--Chinczyna. Jej opis wymaga $\AC G$ i pochodnej aproksymatywnej. Wąska całka Denjoya, równoważna HK i Perronowi, wymaga $\AC G_*$ oraz zwykłej pochodnej prawie wszędzie. Przyjęcie tylko $\AC G$ i zwykłej pochodnej bez dodatkowych warunków prowadzi do innego schematu definicyjnego; obu pojęć nie należy utożsamiać bez dowodu.
\end{uwaga}

